\documentclass[10pt,a4paper]{amsart}
\usepackage[margin=19mm]{geometry}
\usepackage{amsmath,amssymb,mathtools}
\usepackage[scr=rsfs,cal=euler]{mathalfa}
\usepackage{xfrac}
\usepackage[unicode,hidelinks,bookmarks=false]{hyperref}
\newcommand{\ie}{i.e.\ }
\newcommand{\cf}{cf.\ }
\newcommand{\resp}{respectively\ }
\newcommand{\Subsection}{\subsection}
\providecommand{\nobreakdash}{\nobreak\hbox{-}}
\setlength{\emergencystretch}{2em}
\multlinegap0pt
\usepackage{amssymb}
\usepackage[shortlabels]{enumitem}
\usepackage{xcolor}
\newtheorem{lem}{Lemma}
\newtheorem{prop}{Proposition}
\newtheorem{cor}{Corollary}
\def \CC {\mathbb{C}}
\DeclareMathOperator{\Gr}{Gr}
\def \Ical {\mathcal{I}}
\def \Lcal {\mathcal{L}}
\def \Ocal {\mathcal{O}}
\def \PP {\mathbb{P}}
\DeclareMathOperator{\supp}{supp}
\title{The secant variety and syzygies of the Hilbert scheme of two points}
\author{Chiwon Yoon, Haesong Seo}
\date{Source: arXiv:2403.07315v3 (CC BY 4.0)}
\begin{document}
\maketitle

\noindent\emph{Source and licence.} The secant variety and syzygies of the Hilbert scheme of two points. Version arXiv:2403.07315v3 (CC BY 4.0), distributed by arXiv under the Creative Commons Attribution 4.0 International licence, http://creativecommons.org/licenses/by/4.0/, as recorded in the arXiv licence metadata for this exact version and shown on the abstract page https://arxiv.org/abs/2403.07315v3. Full licence notice: \texttt{licenses/arxiv-2403.07315-CC-BY-4.0.txt}.\par\smallskip

\noindent\textit{Editorial scope.} Counted unit: the article's prop prop:Secant-lines-intersect (source0.tex lines 1469--1471). The article's complete original proof of it is delivered with it (source0.tex lines 1473--1568, one \texttt{proof} environment of 96 lines). No mathematics is written, repaired or re-derived: the statement and the proof are the article's own lines, in the article's own notation, and no retained line is substituted or re-flowed.  Retained uncounted as the source-local context the proof needs: lines 1051--1051; lines 1259--1266; lines 1285--1289; lines 1299--1302; lines 1304--1308; lines 1375--1377.  Nothing was omitted from any retained span, so the package carries no omission record.  No retained line contains a citation, so the wrapper carries no bibliography and no bibliography entry.  No retained line contains a figure environment, an image-inclusion command or drawing code, so no image is delivered and the package declares no asset dependency.  Editorial formatting only, in addition: an AMS \texttt{amsart} preamble that restates the article's own declarations and macros used by the retained lines, namely the environments \texttt{lem}, \texttt{prop}, \texttt{cor} on separate counters, together with the standard packages those lines need.  The retained environments print in this document as Lemma 1, Proposition 1, Corollary 1 in order of appearance, whereas the article prints the same items with the numbering of its own full document.  The complete native source of the article as distributed in the arXiv e-print for this exact version is bundled unchanged as \texttt{sources/155941/source0.tex}.
\subsection{Identifiability of the Grassmannian.} \label{subsec:identifiability-of-grassmannians}

\begin{lem} \label{lem:sum-of-ideals}
    Let $\Lcal$ be a $d$-very ample line bundle on $X$ and let $Z,W$ be $0$-cycles on $X$.
    If the colength of $\Ical_Z \cap \Ical_W$ is at most $d+1$, then 
     \begin{equation}
        H^0(\Lcal \otimes (\Ical_{Z}+\Ical_{W})) = H^0(\Lcal \otimes \Ical_{Z}) + H^0(\Lcal \otimes \Ical_{W})
     \end{equation}
    in $H^0(\Lcal)$.
\end{lem}

\begin{lem} \label{lem:Hamming-distance-and-length}
    Let $0 \leq d' \leq d$ be an integer.
    Let $\Lcal$ be a $\min \{ d+d', 2d-1 \}$-very ample line bundle on $X$, and let $Z_1, Z_2 \subset X$ be $0$-cycles of length $d$ whose Hamming distance in $X^{[d]}$ is $d'$.
    Then $\Ical_{Z_1} \cap \Ical_{Z_2}$ has colength $d+d'$ and $\Ical_{Z_1} + \Ical_{Z_2}$ has colength $d-d'$.
\end{lem}

\begin{prop} \label{prop:lines-in-Grassmannian}
    Let $V$ be a vector space of dimension $n$ over $\CC$ and let $1 \leq k \leq \frac{n}{2}$.
    Under the Pl\"ucker embedding $\Gr(k,V) \hookrightarrow \PP^N$, a line $\ell \subset \Gr(k,V)$ corresponds to a pair of subspaces $W_0 \subset W_1 \subset V$ of dimension $k-1$ and $k+1$, respectively, and a point $[U] \in \ell$ corresponds to a $k$-dimensional subspace $W_0 \subset U \subset W_1$.
\end{prop}

\begin{lem} \label{lem:lines-in-Hilbert-scheme}
    Assume that $\Lcal$ is a $(d+1)$-very ample line bundle on $X$.
    For given distinct two points $[Z_1], [Z_2] \in X^{[d]}$, the line $\langle [Z_1],[Z_2] \rangle$ is contained in $X^{[d]}$ if and only if $\mathrm{dist}([Z_1],[Z_2]) = 1$ and $Z_1$ differs from $Z_2$ only at one point of $X$, i.e., $\Ocal_{Z_1, p} = \Ocal_{Z_2, p}$ except for a single point $p \in X$.
    Moreover, if a line meets $X^{[d]}$ in at least three distinct points, then it is contained in $X^{[d]}$.
\end{lem}

\begin{cor} \label{cor:No-lines-for-curve}
    If $X$ is a curve, then no three points in $X^{[d]}$ are colinear.
\end{cor}

\begin{prop} \label{prop:Secant-lines-intersect}
    Two distinct secant lines to $X^{[d]}$ do not intersect out of $X^{[d]}$.
\end{prop}

\begin{proof}
Let $[Z_i] \in X^{[d]}$, $1 \leq i \leq 4$ be four points.
For simplicity, let $\ell_{ij}$ denote the line $\langle [Z_i], [Z_j] \rangle$ and $d_{ij}$ denote the distance $\mathrm{dist}([Z_i], [Z_j])$ for $1 \leq i, j \leq 4$.
Assume that $\ell_{12}$ and $\ell_{34}$ meet outside $X^{[d]}$.
From Section~\ref{subsec:identifiability-of-grassmannians}, we infer that $d_{12} = d_{34} \in \{ 1,2 \}$.

\hypertarget{Case1-1}{\textit{Case 1.}} Assume that 
$d_{12} = d_{34} = 1$.
Let $Z_{12}$ (resp. $Z_{34}$) be the length $d-1$ subscheme defined by $\Ical_{Z_1} + \Ical_{Z_2}$ (resp. $\Ical_{Z_3} + \Ical_{Z_4}$); and let $W_{12}$ (resp. $W_{34}$) be the length $d+1$ subscheme defined by $\Ical_{Z_1} \cap \Ical_{Z_2}$ (resp. $\Ical_{Z_3} \cap \Ical_{Z_4}$).
By Proposition~\ref{prop:lines-in-Grassmannian}, the points $[V] \in \ell_{12}$ correspond to the subspaces $V \subset H^0(\Lcal)$ of codimension $d$ with
 \[
  H^0(\Lcal \otimes \Ical_{W_{12}}) \subset V \subset H^0(\Lcal \otimes \Ical_{Z_{12}}),
 \]
and a similar result holds for $\ell_{34}$.
Let $[V] \in \ell_{12} \cap \ell_{34}$ be the intersection point.
Since the two lines are distinct, we have either $Z_{12} \neq Z_{34}$ or $W_{12} \neq W_{34}$.

If $Z_{12} \neq Z_{34}$, it follows that
 \[
  V = H^0(\Lcal \otimes \Ical_{Z_{12}}) \cap H^0(\Lcal \otimes \Ical_{Z_{34}})
    = H^0(\Lcal \otimes (\Ical_{Z_{12}} \cap \Ical_{Z_{34}})),
 \]
so $[V]$ belongs to $X^{[d]}$.

If $W_{12} \neq W_{34}$, then they have Hamming distance $1$ in $X^{[d+1]}$.
According to Lemma~\ref{lem:Hamming-distance-and-length}, $\Ical_{W_{12}} \cap \Ical_{W_{34}}$ has colength $d+2$.
Then by Lemma~\ref{lem:sum-of-ideals}, 
 \[
  V = H^0(\Lcal \otimes \Ical_{W_{12}}) + H^0(\Lcal \otimes \Ical_{W_{34}}) = H^0(\Lcal \otimes (\Ical_{W_{12}} + \Ical_{W_{34}})).
 \]
Hence we would get a contradiction in any cases.

\hypertarget{Case1-2}{\textit{Case 2.}} Now assume that $d_{12} = d_{34} = 2$.
Write
 \[
  [Z_1] = [v_1 \wedge \cdots \wedge v_{d-2} \wedge v_{d-1} \wedge v_d], \qquad
  [Z_2] = [v_1 \wedge \cdots \wedge v_{d-2} \wedge v_{d+1} \wedge v_{d+2}]
 \]
for some linearly independent vectors $v_1, \dots, v_{d+2} \in H^0(\Lcal)^{\vee}$.
We will apply appropriate linear coordinate changes to $v_i$ as new expressions emerge.
The intersection point $P = \ell_{12} \cap \ell_{34}$ is
 \[
  [v_1 \wedge \cdots \wedge v_{d-2} \wedge (v_{d-1} \wedge v_d + v_{d+1} \wedge v_{d+2})],
 \]
so the annihilator of the kernel of $\wedge P : H^0(\Lcal)^{\vee}\rightarrow \bigwedge^{d+1} H^0(\Lcal)^{\vee}$ is
\[
    H^0(\Lcal \otimes \Ical_{Z_1})+H^0(\Lcal \otimes \Ical_{Z_2})=H^0(\Lcal \otimes \Ical_{Z_3})+H^0(\Lcal \otimes \Ical_{Z_4}).
\]
In particular, $d_{ij} \leq 2$ for any $i,j$.

\hypertarget{Case1-2-a}{\textit{Case 2-a.}} If $d_{13} = d_{23} = 1$, the expression for $[Z_3]$ would be
 \[
  [Z_3] = [v_1 \wedge \cdots \wedge v_{d-2} \wedge v_{d-1} \wedge v_{d+1}].
 \]
Since $[Z_4] \in \langle P, [Z_3] \rangle$, the vector
 \begin{align*}
  &v_1 \wedge \cdots \wedge v_{d-2} \wedge (v_{d-1} \wedge v_d + v_{d+1} \wedge v_{d+2} + av_{d-1} \wedge v_{d+1}) \\
  &= v_1 \wedge \cdots \wedge v_{d-2} \wedge (v_{d-1} \wedge (v_d + av_{d+1}) + v_{d+1} \wedge v_{d+2})
 \end{align*}
should be decomposable for some $0 \neq a \in \CC$.
But since the vectors $v_1, \dots, v_{d-1}$, $v_{d}+av_{d+1}$, $v_{d+1}$ and $v_{d+2}$ are linearly independent, this is not the case.

\hypertarget{Case1-2-b}{\textit{Case 2-b.}} If $d_{13} = 1$ and $d_{23} = 2$, one can write
 \[
  [Z_3] = [v_1 \wedge \cdots \wedge v_{d-2} \wedge v_{d-1} \wedge (v_{d} + v_{d+1})]
 \]
and
 \[
  [Z_4] = [v_1 \wedge \cdots \wedge v_{d-2} \wedge v_{d+1} \wedge (v_{d-1} + v_{d+2})].
 \]
Hence $d_{13} = d_{24} = 1$ and $d_{14} = d_{23} = 2$.
Since the four points $[Z_i]$ are coplanar, $\ell_{13}$ and $\ell_{24}$ must intersect.
They intersect within $X^{[d]}$ by Case~\hyperlink{Case1-1}{1}, and thus they are contained in $X^{[d]}$ by Lemma~\ref{lem:lines-in-Hilbert-scheme}.
In particular, all the $Z_i$ have the same supports.

The case $n = 1$ does not happen by Corollary~\ref{cor:No-lines-for-curve}.

If $d = 2$, we get a contradiction as $\supp(Z_1)$ is disjoint from $\supp(Z_2)$.

\hypertarget{Case1-2-c}{\textit{Case 2-c.}} Finally, suppose that $d_{13} = d_{23} = 2$.
We may assume that $d_{14} = d_{24} = 2$, otherwise we can argue as in Case \hyperlink{Case1-2-a}{2-a} and \hyperlink{Case1-2-b}{2-b}.

If $n=1$, from the equality
 \[
  H^0(\Lcal \otimes (\Ical_{Z_1} \cap \Ical_{Z_2})) = H^0(\Lcal \otimes (\Ical_{Z_3} \cap \Ical_{Z_4})),
 \]
we have $Z_1 \cap Z_2 = Z_1 \cap Z_3$ and $Z_1 \cup Z_2 = Z_1 \cup Z_3$.
Hence we have
 \[
  Z_2 = (Z_1 \cup Z_2) + (Z_1 \cap Z_2) - Z_1 = (Z_1 \cup Z_3) + (Z_1 \cap Z_3) - Z_1 = Z_3
 \]
as divisors on $X$, which is a contradiction.

When $d = 2$, the supports of $Z_i$ are pairwise disjoint for $1 \leq i \leq 4$.
As above, we have $Z_1 \cup Z_2 = Z_3 \cup Z_4$ by $4$-very ampleness, which is impossible.
\end{proof}

\end{document}
